% !TEX root = ../main.tex
\section{Experiments and recasting}\label{sec:exp}
In this section, we present the details of existing HNL searches and classify them according to their search 
strategy. We consider experiments employing the signatures: (i) direct decays, (ii) displaced vertices, and (iii) 
missing energy. For the direct-decay searches, we further partition our analysis into: (i.a) peak searches and (i.b) 
branching ratio searches. The displaced-vertices searches are also split further into: (ii.a) beam-dump searches 
and (ii.b) collider searches.
Missing-energy searches allow us to derive new and stronger constraints on single RPV couplings. In addition, we present 
selected benchmark scenarios for which we will obtain single-coupling and coupling-product bounds within the RPV-MSSM 
framework. We provide an  overview of the experiments, discussed in this section, in Table~\ref{tab:exp_summary}.


\begin{table}[h!t]
\centering
\begin{tabular}{|c|lll|cr|}
\hline
Search Strategy & Ref. & Experiment &  Status & HNL Mixing  & HNL Mass region   \\
 \hline
\multirow{7}{*}{Peak} & \cite{PIENU:2017wbj}& \texttt{PIENU} & curr. & $|U_e|$ &  \SIrange{65}{153}{\MeV} \\ 
&\cite{PIONEER:2022alm} & \texttt{PIONEER} & proj. &  $|U_\mu|$& \SIrange{15.7}{33.8}{\MeV} \\
& \cite{PIONEER:2022alm} & \texttt{PIONEER} & proj. & $|U_e|$ & \SIrange{65}{135}{\MeV} \\
&\cite{Daum:1987bg} & \texttt{SIN} & curr. & $|U_\mu|$ & \SIrange{1}{16}{\MeV}\\
&\cite{NA62:2020mcv} & \texttt{NA62} & curr.   & $|U_\mu|$ & \SIrange{144}{462}{\MeV}\\
&\cite{NA62:2021bji} & \texttt{NA62} & curr.  & $|U_e|$ & \SIrange{200}{384}{\MeV}\\
&\cite{Asano:1981he}& \texttt{KEK} & curr.   & $|U_\mu|$ & \SIrange{160}{230}{\MeV} \\
&\cite{Hayano:1982wu} & \texttt{KEK} & curr.   & $|U_\mu|$ & \SIrange{70}{300}{\MeV}\\
\hline
\multirow{2}{*}{Branching Ratio} &\cite{PIENU:2017wbj} & \texttt{PIENU} &  curr. &$|U_e|$ &  \SIrange{0}{65}{\MeV} \\ 
& \cite{PIONEER:2022alm} & \texttt{PIONEER} & proj.  & $|U_e|$ & \SIrange{0}{65}{\MeV} \\
\hline
\multirow{9}{*}{Beam-dump} & \cite{Ballett:2019bgd} & \texttt{DUNE} & proj. &  $|U_e|$, $|U_\mu|$, $|U_\tau|$ & \SIrange{0}{1968.34}{\MeV} \\
& \cite{T2K:2019jwa} & \texttt{T2K} & curr. &   $|U_e|, |U_\mu|$ & \SIrange{10}{490}{\MeV} \\
& \cite{CHARM:1985nku,CHARMII:1994jjr} & \texttt{CHARM} & curr. &  $|U_e|$, $|U_\mu|$ & \SIrange{300}{1869.65}{\MeV} \\
& \cite{Orloff:2002de} & \texttt{CHARM} & curr. &  $|U_\tau|$ & \SIrange{290}{1600}{\MeV} \\
& \cite{NuTeV:1999kej} & \texttt{NuTeV} & curr. &  $|U_\mu|$ & \SIrange{259}{2000}{\MeV} \\
& \cite{Kelly:2021xbv} & \texttt{MicroBooNE} & curr. &  $|U_\mu|$ & \SIrange{20}{200}{\MeV}\\
& \cite{WA66:1985mfx} & \texttt{BEBC} & curr. &  $|U_e|, |U_\mu|$ & \SIrange{500}{1750}{\MeV}\\
& \cite{Barouki:2022bkt} & \texttt{BEBC} & curr. &  $ |U_\tau|$ & \SIrange{100}{1650}{\MeV}\\
& \cite{Coloma:2019htx} & \texttt{SK} & curr. &  $|U_e|$, $|U_\mu|$ & \SIrange{150}{400}{\MeV} \\
\hline
\multirow{3}{*}{Collider} & \cite{Kling:2018wct} & \texttt{FASER} & proj. & $|U_e|$, $|U_\mu|$, $|U_\tau|$ & \SIrange{0}{6274.9}{\MeV} \\
& \cite{DeVries:2020jbs} & \texttt{MoEDAL-MAPP1} & proj. &  $|U_e|$ & \SIrange{0}{6274.9}{\MeV}\\
& \cite{BaBar:2022cqj} & \texttt{BaBar} & curr. &  $|U_\tau|$ & \SIrange{100}{1360}{\MeV}\\
\hline
\multirow{2}{*}{Missing Energy} & \cite{NA62:2020pwi} & \texttt{NA62} & curr. &   $\text{BR}(\pi^0 \to \text{inv.})$ & \SIrange{0}{134.97}{MeV}\\
& \cite{BaBar:2012yut} & \texttt{BaBar} & curr. &$\text{BR}(B^0 \to \text{inv.})$ & \SIrange{0}{5279.65}{MeV}\\
\hline
\end{tabular}
\caption{Summary of experiments reviewed in Sec. \ref{sec:exp}, sorted by search strategy. We list the relevant references, the status of derived bounds (current or projected), the relevant HNL mixing, and the experimentally accessible HNL mass range.}
\label{tab:exp_summary}
\end{table}




\subsection{Direct-decay searches}

One of the main ways to produce light HNLs is via the decay of light mesons such as pions and kaons.
In direct searches, a beam of charged mesons is brought to a stop inside a scintillator where the mesons decay at rest, or 
the beam mesons are tagged and their positions, momenta, and timing information are measured by a silicon pixel spectrometer.


The energy spectrum of the visible secondary particle, \textit{i.e.}, a muon or an electron, arising from these meson 
decays is measured.  The signal shape of the energy spectrum can be compared with Monte-Carlo simulations for different 
HNL mass hypotheses. Finding no extra peaks in the secondary energy spectrum or rejecting each mass hypothesis allows us to 
exclude the relevant HNL parameters. 


\subsubsection{Peak searches}
%
In peak searches, the energy spectrum of the secondary particle is scanned for additional peaks hinting at HNLs. 

\begin{itemize}
    \item At the 
Swiss Institute for Nuclear Research (\texttt{SIN}), a pion beam was used to put bounds on the mixing $|U_\mu|^2$ in the 
HNL mass range of \SIrange{1}{16}{\MeV}~\cite{Daum:1987bg}. 

\item A search for massive neutrinos at the \texttt{PIENU} experiment~\cite{PIENU:2017wbj} has been made in the decay of pions into positrons.
No evidence was found for additional peaks in the positron energy spectrum.
Thus, upper limits at 90\% confidence level (CL) on $|U_e|^2$ were derived in the HNL mass region \SIrange{60}{135}{\MeV}. 
In another analysis of the \texttt{PIENU} experiment~\cite{PIENU:2019usb}, heavy neutrinos were searched for in pion decays into muons.
The energy spectrum did not show any additional peaks other than the expected peak for a light neutrino.
Thus, the analysis derived a bound on $|U_\mu|^2$ for the HNL mass range of \SIrange{15.7}{33.8}{\MeV}.

\item The \texttt{PIONEER}~\cite{PIONEER:2022alm} experiment is a next-generation rare pion decay experiment. The experiment will perform the 
same search strategy as the \texttt{PIENU} experiment with higher statistics and significantly suppressed background. A 
peak search in the positron spectrum will allow probing $|U_e|^2$ in the HNL mass region \SIrange{65}{135}{\MeV}. 
Further, a search for an additional peak within the muon energy spectrum will allow us to test $|U_\mu|^2$ for $\SI{15.7}
{\MeV} < m_N < \SI{33.8}{\MeV}$.

\item \texttt{KEK}~\cite{Asano:1981he} derived an upper bound on $|U_\mu|^2$ for a massive HNL in the mass range of \SIrange{160}{230}{\MeV}.
A similar search at \texttt{KEK}~\cite{Hayano:1982wu} led to an upper bound on $|U_\mu|^2$ in the HNL mass range of \SIrange{70}{300}{\MeV}.

\item Using a kaon beam, the \texttt{NA62} collaboration placed bounds on $|U_e|^2$ for an HNL with a mass of \SIrange{144}{462}{\MeV}~\cite{NA62:2020mcv}. The analysis approach is different from \texttt{PIENU}.
In this case, a peak-search procedure measures the $K^{+} \rightarrow e^{+} N$ decay rate with respect to the $K^{+} \rightarrow e^{+} \nu$ rate for an assumed HNL mass $m_N$. The HNL mass is varied over the mentioned mass range. The benefit of this approach is the cancellations of residual detector inefficiencies, as well as trigger inefficiencies, and random veto losses.
A similar analysis~\cite{NA62:2021bji} has been performed to measure $|U_\mu|^2$ within the HNL mass range of \SIrange{200}{384}{\MeV}. 
Note that both bounds of \texttt{NA62} are derived with the assumption that the lifetime of the neutral particle exceeds \SI{50}{\ns}.

 



\item The \texttt{BaBar} experiment~\cite{BaBar:2022cqj}  at SLAC has performed a search for the rare decay  
$\tau^- \to \pi^- \pi^- \pi^+ + N$ in the mass region of $100 < m_N < \SI{1360}{\MeV}$. The observed kinematic phase 
space distribution of the hadronic system allows \texttt{BaBar} to place a stringent bound on $|U_\tau|^2$. However, the 
search is based on three-prong tau events. Technically, this allows us to derive a single coupling bound on $\lambda^
\prime_{311}$. However, the four-body production mode of the light neutralino would need proper phase-space 
consideration. Therefore, we do not include a reinterpretation of this search in our work and shall discuss it
elsewhere.

\end{itemize}



\subsubsection{Branching-ratio searches}
%
%
It is possible to measure branching ratios of different pion decay modes. The ratio 
\begin{align}
R_{e/\mu} = \frac{\Gamma\left(\pi^+ \to e^+ + \nu(\gamma)\right)}{\Gamma\left(\pi^+ \to \mu^+ + \nu(\gamma)\right)},
\end{align}
can be used to derive limits on the mixing $|U_e|^2$ in the region $m_N<\SI{65}{\MeV}$.
This has been performed by \texttt{PIENU}~\cite{PiENu:2015seu} and is planned for \texttt{PIONEER}~\cite{PIONEER:2022alm}.




\subsection{Displaced-vertex searches}
Beam-dump and collider experiments can produce HNLs via the same processes that produce light neutrinos. A proton beam 
hitting a fixed target typically produces a large number of pions and kaons, and also heavier mesons. If kinematically 
allowed, the decay of the primary mesons can produce HNLs, which will propagate freely since they are long-lived and 
interact only feebly. The HNLs produced at beam-dump experiments are typically boosted in the forward direction, which 
further increases their decay length in the lab frame. Hence, only a fraction of the produced HNLs decay at the location 
of the detector. To reduce possible background events, the experiments usually have a system of veto detectors equipped 
for both charged and neutral particles. The search strategy relies on the visibility of the HNL decay products inside 
the detector; we will discuss these later.

In the following, we present displaced-vertex searches at beam-dump and collider 
experiments separately.



\subsubsection{Beam-dump search}
%
%
\begin{itemize} 
\item At \texttt{DUNE}~\cite{DUNE:2020lwj}, HNLs can be produced via pion-, kaon-, and  
$D$-meson-decays.
Ref.~\cite{Ballett:2019bgd} predicts constraints on HNLs by searching for their decay products inside the \texttt{DUNE} Near Detector.
It is assumed that only the three-neutrino final state is not detectable.
The search strategy allows to measure all three mixings $|U_e|^2, |U_\mu|^2$ and $|U_\tau|^2$ for an HNL mass range up to the mass of the $D_s$-meson.
Note that in the analysis, a single mixing element is assumed to dominate over the other two at a time.

\item The \texttt{T2K} experiment~\cite{T2K:2019jwa} follows the same approach but uses a kaon beam.
Thus, they derive bounds on both $|U_e|^2$ and $|U_\mu|^2$ for $\SI{10}{\MeV} < m_N < \SI{490}{\MeV}$.


\item Heavy neutral leptons in the \texttt{CHARM} beam-dump experiment are produced from $D$ and 
$D_s$-meson decays.\footnote{In this work, we use ``$D$-mesons'' (``$B$-mesons'') to label the $D^\pm$ 
mesons ($\overset{\scriptscriptstyle(-)}{B^0}$ and $B^\pm$ mesons), while $D_s$ and $B_c$ mesons are
separately discussed. We do not take into account $\overset{\scriptscriptstyle(-)}{D^0}$ or $\overset 
{\scriptscriptstyle(-)}{B_s}$ mesons.} The former allows to set bounds on $|U_e|^2, |U_\mu|^2$ 
for an HNL mass of \SIrange{300}{1869.65}{\MeV}~\cite{CHARM:1985nku,CHARMII:1994jjr}, and the latter 
allows to probe $|U_\tau|^2$ for $\SI{290}{\MeV} < m_N < \SI{1600}{\MeV}$~\cite{Orloff:2002de}.

\item A search for HNLs has been performed at the \texttt{NuTeV} experiment~\cite{NuTeV:1999kej} at 
Fermilab. The data were examined for HNLs decaying into muonic final states to derive bounds on the 
mixing $|U_\mu|^2$ of HNLs in the \SIrange{0.25}{2.0}{\GeV} mass range.
See also Ref.~\cite{Dedes:2001zia} which directly considers the \texttt{NuTeV} data in terms of a light neutralino.


\item An analysis of current data from the \texttt{MicroBooNE} experiment~\cite{Kelly:2021xbv} can 
constrain the parameters of HNLs, that mix predominantly with muon-flavored 
neutrinos, for HNL masses between \SIrange{20}
{200}{\MeV}. 


\item The \texttt{BEBC} experiment derived limits on the HNL-light neutrino mixing parameters from a search for decays of heavy neutrinos in a proton beam-dump 
experiment~\cite{WA66:1985mfx}. It derived bounds on $|U_e|^2, |U_\mu|^2$ for an HNL mass between 
\SI{0.5}{\GeV} and \SI{1.75}{\GeV}. A re-analysis~\cite{Marocco:2020dqu, Barouki:2022bkt} has 
demonstrated that the \texttt{BEBC} detector was also able to place bounds on the $|U_\tau|^2$ 
mixing for HNL masses higher than the kaon mass. This re-analysis has taken into account several 
production and decay channels of HNLs.


\item For \texttt{Super-Kamiokande (SK)}, the largest contribution to HNL production is 
through the decay of mesons produced in the atmosphere via cosmic rays.
A secondary contribution to the flux comes from the HNL production in neutral-current scattering of atmospheric neutrinos passing through the Earth.
The total number of HNL decays inside the detector within a given time window results in an upper bound on the HNL mixing. This approach is different from the displaced-decay search limits from beam dumps and allows one to derive bounds on $|U_e|^2$ and $|U_\mu|^2$ in the minimal HNL scenarios for masses between \SIrange{150}{400}{\MeV}~\cite{Coloma:2019htx}.

\end{itemize}



\subsubsection{Collider searches}
%
%
HNLs can be also searched for in (semi-)leptonic decays of mesons produced at the LHC via a similar mechanism 
as in beam-dump experiments.
At colliders, particle collisions produce an abundance of mesons, such as pions, kaons, $D$-mesons, and $B$-mesons, which can further decay into HNLs.
Far detectors at the LHC are sensitive to decaying, light LLPs with a decay length  comparable
with their distance to the interaction point, \textit{e.g.},  of $\mathcal{O}(1)-\mathcal{O}(100)$ \SI{}{\m}. 
HNLs decaying inside the detector can be identified by their decay products, except for the 
invisible three-neutrino final state.
To ensure a low background environment, the experiments are proposed to be set far away from the primary proton-proton collision points and require shielding between the interaction points and the detectors.

Some of the new detectors at the LHC are already approved and currently running: \texttt{FASER}~\cite{Feng:2017uoz} 
and \texttt{MoEDAL-MAPP1}~\cite{Pinfold:2019nqj}. We include these running HNL searches in order to 
reinterpret their projected sensitivity.
Their follow-up programs, \texttt{FASER2}~\cite{FASER:2018eoc} and \linebreak {\texttt{MoEDAL-MAPP2}~\cite{Pinfold:2019zwp}}, have been proposed for operation during the high-luminosity LHC phase, with an expected final integrated luminosity \SI{3}{\per\atto\barn} and \SI{300}{\per\femto\barn} respectively, and are taken into account in our numerical analysis, as well.


\texttt{FASER(2)}~\cite{Kling:2018wct}\footnote{In principle, Ref.~\cite{FASER:2018eoc} employs the most 
updated geometrical setup of \texttt{FASER(2)}, but the results shown therein for the HNLs do not separate 
the contributions from $D$- and $B$-mesons. Therefore, we have chosen to reinterpret the results given in 
Ref.~\cite{Kling:2018wct} for the HNLs from the heavy mesons' decays, which are only slightly different 
from those given in Ref.~\cite{FASER:2018eoc}.} is intended to detect long-lived 
particles decaying inside the detector volume.
A sensitivity estimate, taking the detector geometry into account, leads to a specific reach in $|U_e|^2, |U_\mu|^2$, and $|U_\tau|^2$ for an HNL with a mass up to the mass of the $B_c$-meson.
Similarly, \texttt{MoEDAL-MAPP1(2)}~\cite{DeVries:2020jbs,Beltran:2023nli} can probe the mixing $|U_e|^2$ for HNL being produced in decays of $D$- and $B$-mesons for the same mass range as \texttt{FASER(2)}.



In addition, other experimental proposals for LLP far detectors include \texttt{MATHUSLA}~\cite{Chou:2016lxi, 
Curtin:2018mvb, MATHUSLA:2020uve}, \texttt{ANUBIS}~\cite{Bauer:2019vqk}, \texttt{CODEX-b}~\cite{Gligorov:2017nwh}, 
and \texttt{FACET}~\cite{Cerci:2021nlb}. None of these proposed detectors has been officially approved. We 
therefore do \textit{not} include them in our reinterpretation strategy. We still want to 
emphasize the prospect of these experiments in the search for light long-lived particles. Sensitivity estimates, 
worked out in detail on the minimal HNL scenarios in Ref.~\cite{DeVries:2020jbs}, would also provide possible 
discovery potential for the discussed light neutralino scenarios, as worked out for example in 
Refs.~\cite{Dercks:2018eua,Dercks:2018wum, Dreiner:2020qbi,Dreiner:2022swd}. These potential future experiments all intend to 
look for various BSM signatures which include neutralino decays induced by all of the $LH_u$, $LL \bar E$, 
and $LQ\Bar{D}$ operators.

\medskip


\medskip

\subsection{Missing-energy searches}

The \texttt{NA62} search~\cite{NA62:2020pwi} allows us to derive bounds on the branching ratio of pions 
decaying into an invisible final state. This is achieved by the reconstruction of the charged particles 
in the process $K^+ \to \pi^+ \pi^0$. The analysis relies on the tracking of the charged $K^+$ and $\pi
^+$ and can probe $\pi^0$ decay to any invisible final state. \texttt{NA62} 
reported a 90\% CL upper limit on $\text{BR}(\pi^0 \to \text{inv.}) < \SI{4.4e-9}
{}$. The search can be recast to derive bounds on RPV couplings which also contribute to the invisible 
decays of the pion. It turns out that the obtained limits are weaker than the currently existing bounds~\cite{Allanach:1999ic, Dercks:2017lfq} and are, therefore, omitted.

The \texttt{BaBar} experiment has also searched for rare decays $B_0 \to \text{inv.}$~\cite{BaBar:2012yut}.
The search relies on the identification of the other neutral $B$-meson, as the ‘tag side’, and thus, can 
measure a purely invisible decay width of the $B$-meson.
The upper limit at the 90\% CL yields $\text{BR}(B_0 \to \text{inv.}) < \SI{2.4e-5}{}$.
We can recast this search into bounds on the production coupling of a light long-lived neutralino.

We could not find existing searches for invisible decays of the other uncharged 
mesons, \textit{i.e.}, $K_S$, $K_L$, and $B_s$. To derive limits on the decay of these mesons into a 
neutralino we consider the uncertainty of the total decay width of the mesons. We assume all visible 
decay modes to be measured within the decay width. The resulting 
uncertainty can be extended to account for additional invisible decays and potential errors in 
measurements. Thus, we use the uncertainty to establish an upper limit on the branching ratio into 
invisible final states, if kinematically allowed. These could contain 
neutrinos, HNLs, or neutralinos. Assuming the latter saturates the width uncertainty, we derive bounds on 
the $LQ\Bar{D}$ couplings. The measured uncertainties can be found in Ref.~\cite{ParticleDataGroup:2022pth}. 


We could not find bounds from direct searches for massive HNLs in the leptonic 
decays $\mu^-\to e^-+\nu  + N$ and $\tau^- \to e^-/\mu^-+ \nu + N$. There are only searches for the muons 
and the $\tau$ leptons to decay to a lighter charged lepton plus active neutrinos or a 
photon~\cite{MEG:2016leq}. Thus, the existing limits on the branching ratio cannot be generalized for a 
massive HNL owing to kinematic assumptions. In this case, we again rely on the uncertainty of the decay widths. This allows us to derive bounds on the $LL
\bar{E}$ couplings. Again, the uncertainties of the decay width are taken from 
Ref.~\cite{ParticleDataGroup:2022pth}.

{\color{white}.

.}

\subsection{Other searches}
For the bilinear coupling scenarios, we use the most relevant and up-to-date constraints on the HNLs existing in the literature, which are summarized in Ref.~\cite{Fernandez-Martinez:2023phj}.
Therefore, we supplement the results from the previously discussed experiments with data from \texttt{TRIUMF}~\cite{Britton:1992xv}, 
\texttt{PSI}~\cite{Daum:1987bg}, \texttt{Borexino}~\cite{Borexino:2013bot,Plestid:2020ssy}, and atmospheric 
neutrinos scattering in the Earth~\cite{Dentler:2018sju}.





\subsection{The recasting procedure}
The procedures for recasting the HNL bounds into limits on RPV scenarios depend on the RPV couplings that 
are switched on, and the search strategy of the experiments; they fall into one of the following three 
categories: 
%
\begin{itemize}
    \item The most straightforward case is for scenarios involving bilinear RPV couplings. As discussed in 
    Sec.~\ref{sec:model}, these couplings lead to mixing between the neutralino and the neutrinos, 
    \textit{cf}.~Eq.~\eqref{eq:nu_bino_bilinear_mixing}. Thus, we directly translate HNL exclusion limits 
    in the mixing vs.~mass plane into bounds in the RPV coupling vs.~neutralino mass plane, using 
    Eq.~\eqref{eq:bilinearmixingmatrix}.
    %
    \item In RPV scenarios involving $LQ\Bar{D}$ or $LL\Bar{E}$ operators, and a (detector-level) stable 
    neutralino produced in the decay of a meson or lepton, missing-energy searches and peak searches can 
    provide sensitivity. We use the HNL exclusion limits (typically in the mixing vs.~mass plane) in order 
    to determine the bounds on the decay width of the relevant meson/lepton into an HNL (see, for instance, 
    Refs.~\cite{Gorbunov:2007ak,Gorbunov:2020rjx}). Since replacing the HNL with a light bino of the same 
    mass does not change the experimental signature and the kinematics, we can simply equate the above with 
    the corresponding decay width of the same meson/lepton into a bino in the RPV model, using the 
    expressions given in Sec.~\ref{sec:model}; this gives us the bound on the relevant RPV coupling in 
    terms of the bino mass.
    %
    \item For RPV scenarios involving $LQ\Bar{D}$ or $LL\Bar{E}$ operators where the neutralino is produced
    in meson/lepton decays, and is \textit{unstable} at the detector scales, displaced-vertex searches (beam dump or collider) can apply. In this case, we use the long-lifetime approximation to calculate the number 
    of decay events reconstructed in the detector, following the arguments outlined in 
    Ref.~\cite{Beltran:2023nli} (see also Ref.~\cite{Fernandez-Martinez:2023phj}); we review the procedure briefly now. Consider a beam-dump or collider 
    experiment that searches for HNLs via displaced vertices. Let the HNL, $N$, be produced in the decay of 
    a parent particle $P$ (for instance, a pion).\footnote{We are assuming that the direct production of 
    the HNL (or bino) is suppressed compared to indirect production via decays of mesons and leptons, which is 
    typically the case.} Then, the number of detected events for a final state, $Y$, produced in the decay of $N$ can be estimated as,
    \begin{equation}
        N^{\text{HNL}}_{\text{events}} = N_P\times \text{BR}(P\to N+X) \times \text{BR}(N\to Y) \times \epsilon\,,
        \label{eq:HNL_events}
    \end{equation}
    where $N_P$ is the number of $P$ produced at the experiment, $X$ are any additional objects produced in 
    $P$'s decay that are not of interest to us, and the BRs are the corresponding branching ratios. $\epsilon$ is a factor accounting for the detector acceptance and efficiency. This 
    factor is linearly proportional to the probability for the HNL to decay inside the detector ($P[
    \text{decay}]$). It also depends on experiment-specific information such as the detector type and its geometry.
    In the limit of a long-lived HNL, such that the boosted decay length is much 
    larger than the distance $\Delta L$ from the interaction point of the experiment to the first edge of the 
    detector along the direction of travel of the HNL, to a good approximation we have~\cite{Beltran:2023nli},
    \begin{equation}
        P[\text{decay}] \approx \frac{\Delta L \times \Gamma_N}{\beta_N\gamma_N}\,,
    \end{equation}
    where $\Gamma_N, \beta_N$, and $\gamma_N$ represent the total decay width, the relativistic velocity, 
    and the Lorentz boost factor of the HNL, respectively. Inserting this in Eq.~\eqref{eq:HNL_events}, 
    we obtain,
    \begin{equation}
        N^{\text{HNL}}_{\text{events}} = A\times\frac{\text{BR}(P\to N+X) \times \Gamma(N\to Y)}{\beta_N\gamma_N}\,,
        \label{eq:HNL_events_simpler}
    \end{equation}
    where all the quantities depending on the HNL model details are written explicitly, and the remaining 
    experiment-specific factors of Eq.~\eqref{eq:HNL_events_simpler} have been absorbed into the 
    proportionality factor $A$. Now, assume that there is also an RPV scenario in which the light bino, $\tilde
    {\chi}^0_1$, is produced in $P$ decays, and has a decay mode into $Y$. Analogously, we can 
    write,
    \begin{equation}
        N^{\text{RPV}}_{\text{events}} = A\times\frac{\text{BR}(P\to \tilde{\chi}^0_1+X') \times \Gamma(\tilde{\chi}_1^0\to Y)}{\beta_{\tilde{\chi}_1^0}\gamma_{\tilde{\chi}_1^0}}\,,
        \label{eq:RPV_events}
    \end{equation}
    where the HNL is now replaced by the bino. Here, since both the light bino and the HNL are produced in the same meson's decay at the identical experiment, the proportionality factor $A$ can 
    be legitimately assumed to be essentially the same. Thus, combining the two equations, we obtain a simple 
    scaling relation between the two models:
    \begin{equation}
        \frac{N^{\text{RPV}}_{\text{events}}}{N^{\text{HNL}}_{\text{events}}} = \frac{\text{BR}(P\to \tilde{\chi}^0_1+X') \times \Gamma(\tilde{\chi}^0_1\to Y) \times \beta_N\gamma_N}{\text{BR}(P\to N+X) \times \Gamma(N\to Y) \times \beta_{\tilde{\chi}^0_1}\gamma_{\tilde{\chi}^0_1}}\,.
        \label{eq:HNL_RPV_events}
    \end{equation}
    For an HNL and a light bino with the same mass, we can further simplify the above expression in the limit where the HNL and the bino carry the same energy in the lab frame.
    In this case, the $\beta \gamma$ factors cancel out, and we have,
    \begin{equation}
        \frac{N^{\text{RPV}}_{\text{events}}}{N^{\text{HNL}}_{\text{events}}} = \frac{\text{BR}(P\to \tilde{\chi}^0_1+X') \times \Gamma(\tilde{\chi}^0_1\to Y)}{\text{BR}(P\to N+X) \times \Gamma(N\to Y)}\,,
        \label{eq:HNL_RPV_events_simpler}
    \end{equation}
    which is completely free from any experiment-specific factors. This will be the master expression we use for recasting DV searches.
    We note that, technically, having different $X$ and $X'$ induces different kinematics for the HNL and the bino.
    For instance, the two final states could have different masses, or even contain differing number of objects, thus affecting the energy carried by the bino relative to the HNL.
    However, typically, the experiments we consider produce the parent particles (or their decay products) with significant boosts, in which case the above formula is only modified mildly.
    On the other hand, we stress that the HNL and the bino must be produced in the decay of the same (or similar) parents, and must decay into the same final state since these can significantly alter the detector acceptances and efficiencies.\footnote{Actually, the final states into which the HNL and bino decay need not be identical, \textit{e.g.}, the invisible objects contained in both do not need to match; the crucial part is that the detection efficiencies at the considered experiment must be the same. } 
    
    To use Eq.~\eqref{eq:HNL_RPV_events_simpler}, assume that a given DV search for HNLs concludes without discovery, and obtains bounds on the HNL model parameters corresponding to a certain signal-event number, $[N_{\text{events}}^{\text{HNL}}]_{\text{bound}}$.
    Since the signal (and kinematics) in the RPV model is the same, this bound also applies to the light binos: $[N_{\text{events}}^{\text{RPV}}]_{\text{bound}} = [N_{\text{events}}^{\text{HNL}}]_{\text{bound}}$. Plugging this into Eq.~\eqref{eq:HNL_RPV_events_simpler}, we arrive at,
    \begin{align}
        [\text{BR}(P\to \tilde{\chi}^0_1+X') \times \Gamma(\tilde{\chi}^0_1\to Y)]_{\text{bound}}=[\text{BR}(P\to N+X) \times \Gamma(N\to Y)]_{\text{bound}}\,.
    \end{align}
    The right-hand side of the above equation can be evaluated by using the bounds on the HNL mass and mixings as input (see, for instance, Refs.~\cite{Shrock:1980ct, Gorbunov:2007ak, Bondarenko:2018ptm} for explicit expressions for all relevant decay widths in terms of these parameters), while the expressions for the RPV counterparts appear in Sec.~\ref{sec:model} and depend on the RPV couplings, sfermion masses, as well as the neutralino mass.
\end{itemize}
